\section{What the Router Costs}
\label{sec:ch10-what-the-router-costs}
\index{router!latency cost|(}%

Chapter~\ref{ch:the-first-cut-extracting-catalog} split a service in two and
left the price of that unmeasured, on the grounds that half the cost is a hop
nothing was making yet. The hop exists now, so here is the bill.

Five rows, all timed in one run, sixty requests each after fifteen warm-ups.
This is \code{scripts/measure-router.mjs} at tag \code{ch10}, and it is not part
of the verification gate: these are single-machine numbers and asserting them
would make the gate fail on a busier laptop. Medians in milliseconds, printed
twice for a reason that comes after the table:

\begin{minted}[fontsize=\footnotesize]{text}
                                             pass 1   pass 2
catalog direct, titles only                     1.1      1.4
router, titles only (one fetch)                 3.3      4.5
mosaic direct, _entities for 10 products        1.9      2.4
the two calls by hand, in sequence              3.7      3.9
router, the storefront query (two fetches)      6.2      6.4
\end{minted}

\index{measurement!measuring both sides in one sitting}%
Every row is measured in the same sitting, which is
chapter~\ref{ch:data-without-the-n-plus-1}'s rule and not a formality: that
chapter recorded measuring the same code at 10 to 12 ms one afternoon and 19 to
20 ms another, purely because the machine was busier. Rows compared across
sittings say nothing.

The two subtractions the script prints:

\begin{minted}[fontsize=\footnotesize]{text}
router overhead on a single-subgraph query      2.2      3.1
router overhead on the two-hop query            2.5      2.5
\end{minted}

\index{measurement!what survives repetition}%
So the claim that survives is that the router adds roughly two to three
milliseconds to a query on this machine.

The claim that does not survive is anything about one of those rows being
larger than the other. Pass 1 has the two-hop query paying slightly more
overhead, pass 2 has it paying slightly less, and the run-to-run spread on the
first row alone is nearly a millisecond. That is why both passes are printed.
One pass would have supported a tidy story about the cost of fan-out, and it
would have been a story about this machine's scheduler.

\index{measurement!upper bounds}%
All five rows are upper bounds on what a deployment would pay, for the reason
chapter~\ref{ch:how-a-federated-query-actually-runs} recorded about its own
table: the router is in a container reaching the subgraphs across Docker's
bridge and back in through a published port, while the direct rows are loopback
on the host. A router sitting beside its subgraphs pays less than this.

\subsection*{The row that is really the comparison}

\index{router!against assembling by hand}%
The interesting row is the fourth, not the second. \enquote{The two calls by
hand} is what I have been doing since
chapter~\ref{ch:the-first-cut-extracting-catalog}: ask Catalog for the products
and their identifiers, take the identifiers, ask Mosaic for everything else
through \code{\_entities}. It is the same two round trips the router makes, and
it costs 3.7 and 3.9 ms.

Against 6.2 and 6.4 for the router doing it. So the router is not buying
throughput. It is charging a couple of milliseconds to do something a client
could do itself in the same number of network calls.

What it buys is that the client does not know there are two of anything. That is
worth a great deal and it is worth exactly nothing on a latency graph, which is
the honest way to hold both facts. If a team asks what the gateway costs, the
answer here is single-digit milliseconds and a hop they were making anyway; if
they ask what it saves, the answer is not in this table.

\subsection*{The half of the query nobody is measuring}

\index{observability!half a query}%
One more number, and it is the absence of one.
Section~\ref{sec:ch10-the-query-that-stopped-working} showed Mosaic reporting
nineteen resolvers and six SQL statements for its part of the storefront query.
Catalog reports nothing, because chapter~\ref{ch:the-life-of-a-request}'s
instrumentation stayed in \code{Mosaic.Api} when Catalog left.

So the two subtractions above are the only view I have of where the time goes,
and they are measured from outside all three processes. I cannot tell you which
of the 6.2 ms was Catalog, which was Mosaic, which was the router's planning and
which was Docker's bridge, and no amount of running this table again will tell
me. That is what distributed tracing is for, it is
chapter~\ref{ch:observability}'s subject, and the reason to name the gap here is
that a two-service system is exactly where you stop being able to reason about
latency by subtraction and start needing spans.

\index{router!latency cost|)}%
